\chapter{总结与展望}
\label{chap:conclusion}

\section{研究总结}
\label{sec:summary}

本论文针对现有因果发现方法在因果函数建模方面的不足，围绕“如何从观测数据中显式学习因果函数”这一核心问题展开研究。主要工作包括两个方面：

第一，提出了基于 KAN 的连续优化因果发现方法（DAG-KAN）。该方法采用轻量级单层 KAN 架构，通过 B 样条基函数显式拟合变量间的非线性因果机制，同时利用稀疏正则化学习因果邻接矩阵。在模拟和真实数据集上的实验结果表明，在加性模型数据上，DAG-KAN 的结构学习精度优于主流方法，特别是在高维场景下结构汉明距离平均降低了 22.6\%。

第二，提出了基于 KAN 的知识融合因果发现方法（Causal-KAN）。针对真实场景中普遍存在的非参数或半参数数据，该方法通过在输入层引入因果掩码机制，将已知的因果父节点关系编码为网络约束，构建了知识嵌入层与函数拟合层的解耦架构。在模拟和真实数据集上的实验表明，嵌入因果知识能够有效防止过拟合，且学习到的因果函数具有可解释性。

\section{研究局限性}
\label{sec:limitations}

尽管本研究取得了一定进展，但仍存在以下局限性，需要在后续工作中加以改进：

\begin{compactenum}
\item 对先验知识的依赖程度较高
	
Causal-KAN 方法的有效性建立在部分因果结构已知的前提之上。在实际应用中，这种先验知识往往来源于领域专家的经验判断或已有文献的总结，其获取成本较高且可能存在主观偏差。当面对完全未知的因果图谱时，该方法的优势难以发挥。此外，如果先验知识本身存在错误，模型可能将这些错误固化到网络结构中，导致系统性偏差。如何降低对先验知识的依赖，或建立先验知识的可信度评估机制，是亟待解决的问题。
	
\item 高维场景下的计算效率仍有待提升
	
虽然 DAG-KAN 采用了轻量级设计，但在处理超大规模变量（如$d>1000$）或超高维数据（如图像像素级变量）时，计算成本仍然较高。这主要源于两个方面：一是 KAN 的 B 样条基函数需要为每条边维护独立的参数，当变量数量增加时，参数量呈二次方增长；二是连续优化框架需要反复计算梯度并更新参数，在大规模场景下收敛速度较慢。相比之下，纯拓扑发现方法（如 PC 算法）在结构学习阶段的计算效率更高。因此，如何在保持函数建模能力的同时提升计算效率，仍需进一步探索。
	
\item 真实数据中因果函数的评估缺乏客观标准
	
在合成数据实验中，可以通过 MSE、$R^2$ 等指标量化评估因果函数的学习精度，因为真实函数形式是已知的。然而，在 Sachs 等真实数据集中，变量间的真实因果函数是未知的，只能通过生物学理论或专家经验进行定性判断，缺乏定量评估手段。这导致难以客观衡量模型学习到的函数是否准确反映了真实的作用机制，也难以在不同方法之间进行公平比较。
	
	\item 对非加性模型的适用性有限
	
本研究提出的方法主要针对加性噪声模型（ANM）设计，其核心假设是因果机制可以分解为父节点的加性组合。然而，真实世界中存在大量非加性因果机制，如变量间的乘积交互、条件依赖等复杂形式。虽然 KAN 理论上可以逼近任意连续函数，但在非加性场景下，模型的结构设计和正则化策略可能需要重新调整。目前，本研究尚未系统探讨方法在非加性模型上的表现，这限制了方法的适用范围。
\end{compactenum}

\section{未来工作展望}
\label{sec:future_work}

基于上述局限性，未来工作可以从以下几个方面展开：

\begin{compactenum}
	\item 混合框架的构建
	
将 Causal-KAN 与约束型因果发现方法（如 PC 算法、FCI 算法）相结合，形成“从零发现→知识增强”的两阶段混合框架。第一阶段利用约束型方法从数据中初步学习因果骨架，识别可能的因果边；第二阶段利用 Causal-KAN 在已知骨架的基础上精细化学习因果函数。这种混合框架既降低了对先验知识的依赖，又保留了函数建模的优势。
	
	\item 模型效率的优化
	
探索参数共享、低秩分解、稀疏激活等技术，降低 KAN 的参数规模和计算开销。例如，可以研究基于注意力机制的动态稀疏连接策略，仅对重要的因果路径分配计算资源。
	
	\item 评估方法的完善
	
结合领域知识和统计检验，建立真实数据中因果函数的半定量评估标准。例如，可以利用干预实验数据验证学习到的函数是否符合因果效应的方向和量级；或者利用跨环境泛化能力评估函数的稳健性。
\end{compactenum}

综上所述，本研究在因果函数显式建模方面进行了初步探索，但距离构建完善的因果发现与函数学习框架仍有较大差距。希望后续研究能够在上述方向取得突破，推动因果学习从理论走向实践。